\paragraph{Platform}
All measurements use a Tesla T4 (16~GB GDDR6, 40 streaming multiprocessors, PCIe Gen3 $\times$16) on a Google Colab virtual machine with 13.6~GB of host memory, PyTorch \nTorchVersion{} and CUDA \nCudaVersion. Page-locked host-to-GPU copies of a layer (201~MB) run at \nHwBw~GB/s. The T4 has a 70~W power limit, and under sustained load its driver lowers the SM clock to stay within it. In a first set of runs this power capping changed the compute time of a layer by up to \nCapDriftPct\% within one configuration, and a concurrent copy, which itself draws power, slowed computation by 2--9\%. Because such clock changes are a property of this power-limited card rather than of the transfer pipeline, the measurements reported in Section~\ref{sec:results} lock the SM clock with \texttt{nvidia-smi -lgc} at a frequency at which the card stays below its power limit (570~MHz for iterations of up to 1\,024 tokens and 450~MHz for longer ones). A sampler records the SM clock, power and throttle state every 200~ms, so that every configuration can be checked for clock changes. Section~\ref{sec:res-power} reports what happens without the lock.

\paragraph{Model and workloads}
We use the layer shapes of Llama-3.2-3B \cite{dubey2024llama3} (28 layers, model width 3\,072, 24 query and 8 key--value heads of dimension 128, feed-forward width 8\,192) with random FP16 weights; execution time does not depend on the weight values. A unit is one transformer layer: RMS normalisation, the fused query--key--value projection, attention, the output projection, a second normalisation and the gated feed-forward block. Rotary embeddings are omitted. Attention uses the memory-efficient kernel of PyTorch with the key--value heads expanded to the query heads, because the T4 does not support the fused grouped-query path. Prefill iterations process one prompt of $T\in\{256,512,768,1024,1536,2048\}$ tokens with causal attention; decode iterations process one new token for each of $B\in\{16,64\}$ sequences against a key--value cache of 512 tokens that stays in GPU memory. One cache tensor is shared by all layers, which does not change the work per layer and keeps GPU memory free for resident layers. With these shapes the ratio $f/c$ of copy to compute time per layer ranges from \nRatioMin{} to \nRatioMax, which covers the three regimes of Section~\ref{sec:theory}.

\paragraph{Runtime}
The runtime implements assumptions A1--A3 directly. Every layer's weights are one contiguous slice of a single page-locked host buffer. Resident layers are copied to GPU memory once. Streamed layers use $k$ GPU slots of one layer each, a copy stream and a compute stream: copy $j$ waits on an event recorded after the computation of the layer that used slot $j \bmod k$ before it, and the computation of a streamed layer waits on the event recorded after its copy. All work is enqueued asynchronously, so the GPU, not the host, decides when each operation starts; the host enqueue time per iteration is recorded and was always shorter than the iteration. For each configuration the runtime executes two untimed iterations and then 12 timed prefill iterations (40 for decode), discards the first 4 (10), and reports the mean iteration time over the rest. One further iteration records the start and end of every computation and copy with CUDA events, which gives the timelines in Figure~\ref{fig:timeline}.

\paragraph{Configurations}
For every workload we run resident counts $r\in\{0,4,7,10,14,18,21,24\}$, slots $k\in\{2,3\}$ and four placements: the first $r$ layers resident (\emph{contiguous-first}, the order in which model-order filling assigns layers), the last $r$ layers (\emph{contiguous-last}), the balanced placement of Section~\ref{sec:theory} (resident positions $\lfloor (j+\frac12)L/r\rfloor$), and one random placement. A slot sweep at $T=768$ and $r=14$ uses $k\in\{1,2,3,4,6\}$. In total the measurements cover \nAccN{} configurations.

\paragraph{Predictions}
Before the streaming runs of a workload we measure, with all weights resident, the compute time $c_i$ of every layer (median of 7 timed repetitions after warm-up), and once per model the copy time $f_i$ of every layer. The prediction for a configuration is the iteration time given by Lemma~\ref{lem:cycle} for these measured $c_i$ and $f_i$; nothing is fitted to the streaming runs. We also measure every layer's compute time while copies run back to back on the copy stream ($c_i'$) and evaluate a variant of the model in which a layer computes at rate $1/c_i'$ while a copy is in progress and at rate $1/c_i$ otherwise. With locked clocks $c_i'$ and $c_i$ differ by less than \nConcInflationLockedPct\%, so the two predictions coincide; the variant matters only for the unlocked runs.
